% ============================================================
% §2 数据源溯源与精细化特征工程
% 竞赛: 电工杯 B题 (2026)
% 阶段: Stage 01 产物, L1 rubric 待审
% ============================================================

\section{数据源溯源与精细化特征工程}

\subsection{数据资产全景溯源}

本文所依托的原始数据来源于题目附件1至附件5，涵盖某城市成熟街道下辖10个连片小区（编号A-J）的5类异构数据资产。该区域常住总人口31~300人，其中60岁以上老龄人口6~864人，老龄化率达21.9\%，已深度迈入中度老龄化社会。数据结构呈现典型的“微观截面+多维属性”特征——以小区为观测单元，横向展开人口结构、服务需求、空间距离、经济约束与满意度规则5个维度，纵向贯穿自理、半失能、失能3类老年群体，构成$10 \times 3 \times 6$的三维需求张量空间（小区$\times$老人类型$\times$服务项目），为后续马尔可夫演化预测与混合整数选址定价优化提供了坚实的特征基座。

原始数据包含5个Excel工作簿、9个数据表，累计记录210条，采样时间为截面数据（题目假定五年内人口结构与失能率相对平稳变化），缺失率0.0\%。各附件的数据谱系与关键统计摘要如表\ref{tab:data_inventory}所示。

\begin{table}[H]
\centering
\small
\caption{附件数据资产清单与描述性统计摘要}
\label{tab:data_inventory}
\begin{tabular}{p{1.2cm}p{3.5cm}p{1.5cm}p{5.5cm}}
\toprule
\textbf{附件} & \textbf{数据表} & \textbf{维度} & \textbf{关键统计量} \\
\midrule
附件1 & 人口与老人结构 & $10 \times 7$ & 老龄化率21.9\%, 人均月收入均值3250元 \\
       & 转移概率 & $2 \times 1$ & 自理$\to$半失能4.5\%, 半失能$\to$失能10\% \\
\midrule
附件2 & 月均服务需求次数 & $6 \times 3$ & 助餐需求最大(14-22次/月), 紧急救助最小(0.15-3次/月) \\
       & 服务营收与支出 & $6 \times 2$ & 上门护理单价最高(30元), 紧急救助公益免费 \\
       & 月服务消费上限 & $3 \times 1$ & 自理$\leq$20\%, 半失能$\leq$25\%, 失能$\leq$30\%月收入 \\
\midrule
附件3 & 建设与运营成本 & $3 \times 3$ & 小/中/大型: 18/32/45万元建设, 2000/3200/4400元日管理 \\
\midrule
附件4 & 小区间距离矩阵 & $10 \times 10$ & 距离IQR [300, 1800]m, 中位数700m, 77.8\%边$\leq$1000m \\
\midrule
附件5 & 满意度评分规则 & 3因素$\times$4-5档 & $S=0.2S_1+0.3S_2+0.5S_3$, 值域$[0.6, 1.0]$ \\
\bottomrule
\end{tabular}
\end{table}

\subsection{人口结构空间异质性刻画}

10个小区的老人结构并非均匀分布。图\ref{fig:q1_stacked}展示了各小区60岁以上老人的三类状态构成。小区C以920人老龄人口居首，小区F以472人居末，极差达448人——这意味着在后续选址优化中，服务站若优先部署于C、G、E等高密度小区，其服务半径内覆盖的潜在需求基数将显著优于边缘小区。

从失能结构看，全区域自理老人占比68.6\%（4~712人），半失能老人占22.3\%（1~528人），失能老人占9.1\%（624人）。该比例结构与《中国老龄事业发展报告》中城市社区老年人健康状态的宏观分布（自理约65-70\%）高度吻合，侧面验证了附件数据的内部一致性。值得注意的是，半失能与失能老人合计占比达31.3\%——这意味着近三分之一的老年人口对上门护理、康复理疗等高频刚需服务存在强制性依赖，构成了养老服务站运营的核心需求基本盘。

\subsection{服务需求结构的经济约束预诊断}

每位老人月均服务需求次数呈现显著的梯度特征：自理老人月均全量服务消费356元，半失能老人822元，失能老人1~208元。失能老人的月均消费是自理老人的3.4倍，而上限约束（月收入$\times 30\%$）对该群体形成了硬性预算约束。

消费约束预诊断（以区域人均月收入均值3~250元为基准）揭示了一个关键机制性困局：

\begin{enumerate}[label=(\arabic*)]
    \item \textbf{自理老人}: 月预算650元 $>$ 全量消费356元（预算充足率54.8\%），消费约束不起作用，实际需求等于理论需求；
    \item \textbf{半失能老人}: 月预算813元 $\approx$ 全量消费822元（偏差仅1.1\%），处于临界状态——其中低收小区（F区月收入2~700元、D区2~900元）已触发等比例削减机制；
    \item \textbf{失能老人}: 月预算975元 $<$ 全量消费1~208元（缺口23.9\%），在所有10个小区均触发消费约束的等比例削减，其中最低收入小区F的缺口高达49.1\%。
\end{enumerate}

这一”因贫失能”的底层逻辑——失能程度越深、服务需求越大，但消费能力约束越紧——构成了本问题Q1.3需求预测与Q3定价补贴优化的核心经济学张力。后续建模中，我们将严格依据题目给定的”若理论服务费用超消费上限则等比例削减各类服务次数（取整）”规则，对超限群体进行需求量缩放。

\subsection{空间距离矩阵的可达性检验}

附件4给出的小区间距离矩阵为$10 \times 10$对称方阵（对角线为0），经检验矩阵满足严格对称性（$d_{ij} = d_{ji}$, $\forall i,j$），但存在4处轻微的三角不等式违反（例如A$\to$E直达1~500m $>$ A$\to$J$\to$E中转1~300m），这表明部分小区间的实际路径并非欧氏直线距离，可能受到社区道路网络的拓扑约束。鉴于违反幅度均小于200m且仅影响4/360条路径，本文将其视为道路网络的真实迂回效应，不予“修正”。

以服务半径1~000m为可达性阈值，全矩阵90条邻接边中有70条满足可达条件（77.8\%），20条超出服务半径。其中小区A和F各仅有4个可达邻居（可达度最低），而C、D、G、H、I、J均有8个（可达度最高）。这一空间异质性意味着：若将服务站部署于A区或F区，其有效服务覆盖将严重受限；而部署于C区或J区则可最大化服务辐射范围。此发现直接影响了后续Q2 MILP模型的候选站点优先序。

\subsection{满意度评分规则的离散化数学抽象}

题目定义的服务满意度$S = 0.2S_1 + 0.3S_2 + 0.5S_3$为三个因素加权和，值域限定为$[0.6, 1.0]$。三个子满意度均为阶梯状分段常数函数：

\begin{itemize}
    \item \textbf{距离满意度$S_1$}：4档阶梯，断点\{300, 500, 650, 1000\}m，分值\{1.00, 0.90, 0.75, 0.60\}，单调递减；
    \item \textbf{响应满意度$S_2$}：5档阶梯，断点\{0.60, 0.75, 0.85, 0.95, 1.00\}（服务站利用率），分值\{1.00, 0.93, 0.85, 0.72, 0.60\}，单调递减；
    \item \textbf{价格满意度$S_3$}：4档阶梯，断点\{平价, 微溢价10\%, 中溢价20\%, 高溢价\}相对基准价，分值\{1.00, 0.90, 0.75, 0.60\}，单调递减。
\end{itemize}

此规则的离散阶梯特性对后续优化建模提出了严苛的数学要求：不可将其简单近似为连续函数（如距离的倒数或负指数），否则将丧失题目数据的结构化细节。本文在Q2和Q3的MILP建模中，采用Big-M法结合特殊有序集约束（SOS2）对上述分段常数函数进行了精确的线性化编码——将4档/5档/4档的离散跳跃转化为等价的混合整数线性不等式组，使得GLPK求解器可在全局范围内寻优，而非陷入局部近似误差。

\subsection{数据预处理结论}

经上述全量数据审计，形成以下预处理结论：

\begin{enumerate}[label=(\arabic*)]
    \item \textbf{数据完整性}: 5个附件9张表的210条记录均无缺失值，无需填补；
    \item \textbf{数据一致性}: 10个小区三类老人数之和均精确等于60+老人总数，通过交叉校验；
    \item \textbf{空间数据质量}: 距离矩阵对称，4处轻微三角不等式违反保留为道路网络迂回效应；
    \item \textbf{消费约束激活}: 失能老人在全10个小区均触发等比例需求削减，半失能老人在2个低收小区触发，自理老人均不触发——此信息写入后续Q1.3的需求预测模型输入；
    \item \textbf{满意度规则编码}: 确认3因素$\times$(4+5+4)档的离散阶梯结构，后续MILP建模阶段以Big-M+SOS2精确编码，禁用连续近似。
\end{enumerate}
